\chapter{Uso de inteligencia artificial}
\label{cap:uso-ia}

Durante el desarrollo de este Trabajo de Fin de Grado se utilizaron
herramientas de inteligencia artificial generativa como apoyo al desarrollo
de software, al análisis del código existente y a la redacción de esta
memoria. Este capítulo detalla en qué medida y de qué forma se
han utilizado, siguiendo el mismo criterio de transparencia aplicado al
resto del documento. El uso ha sido sobre distintas áreas del proyecto. Las 
decisiones de diseño y de desarrollo han sido siempre del autor y del tutor
mediante su orientación. Sin embargo, las herramientas de IA han sido un apoyo en la redacción de código y 
resolución de dudas de implementación, así como en la redacción de esta memoria.

\section{Herramientas}
\label{sec:ia-herramientas}

La herramienta empleada ha sido \textbf{Claude Code}.

Conviene distinguir dos usos claramente diferenciados a lo largo del
proyecto, que se detallan en las secciones siguientes:

\begin{itemize}
  \item Como asistente de desarrollo de software, en la
        construcción de la aplicación web y revisor de scripts del pipeline.
  \item Como revisor y en casos redactor partes de esta memoria, a partir de una base del autor y de las notas
         que el autor ha tomado durante el desarrollo del proyecto y del análisis directo del código del
        proyecto.
\end{itemize}

No se ha usado ninguna otra herramienta de IA generativa en ningún punto del proyecto.

\section{Uso en el desarrollo de software}
\label{sec:ia-desarrollo}

El pipeline offline en Python es trabajo previo del autor, desarrollado de
forma independiente antes de emplear Claude Code en este proyecto. En esa
parte, el papel de la herramienta se ha limitado a eer, analizar y corregir o editar
el código ya existente para poder documentarlo con precisión en esta
memoria (sección~\ref{sec:ia-memoria}), no a escribirlo.

El uso de Claude Code como herramienta de desarrollo activo se concentra en
la aplicación web (sección~\ref{sec:implementacion-web}), construida 
durante este proyecto con su ayuda:

\begin{itemize}
  \item El paquete \texttt{src/web\_export/} --\texttt{compute\_metadata.py}
        y \texttt{export\_bundle.py}-- se escribió con Claude Code a partir
        de una instrucción explícita: reutilizar la lógica ya validada de
        \texttt{src/metadata/} sin modificar ni duplicar esos scripts,
        parametrizándola por algoritmo. La verificación de que la lógica
        reproducida era efectivamente equivalente se hizo
        contrastando el código generado contra los scripts originales.
  \item Todo el código de base de \texttt{frontend/} se generó igualmente con Claude Code, a partir de
        decisiones de diseño explícitas del autor: el \emph{stack}
        tecnológico (React, Vite, TypeScript), el alcance del MVP fiel al
        boceto ya planificado en el capítulo~\ref{cap:diseno}, y una ronda
        posterior de ajustes visuales concretos (color de acento, cabecera
        con logotipo, alineación de tablas) solicitados de forma iterativa
        tras revisar la aplicación en ejecución.
  \item Antes de implementar la aplicación web se detectó, durante la fase
        de planificación con la herramienta, un desajuste entre el diseño
        original y los datos realmente disponibles. Se plantearon
        al autor dos alternativas.
  \item Durante la implementación se encontraron y corrigieron con la
        herramienta algunos problemas de implementación, como la incompatibilidad del
        \emph{fourcc} \texttt{mp4v} de OpenCV con la reproducción de vídeo
        en navegador, ya documentada como P17 en la
        Tabla~\ref{tab:impl-problemas}.
\end{itemize}

En todos los casos, el código generado se ha ejecutado y verificado
antes de darlo por válido, no
se ha incorporado código sin comprobar que efectivamente hace lo que se le
pidió.

\section{Uso en documentación}
\label{sec:ia-memoria}

El uso de Claude Code en este proyecto ha sido en gran medida en la redacción y revisión de
esta propia memoria.

La herramienta ayudó en consultar ideas y alternativas para la estructura, la redacción de secciones y subsecciones, la revisión de estilo y
la corrección de errores tipográficos y de sintaxis o reestructurar texto a formato LaTeX.
En ningún caso se ha usado para generar contenido técnico o resultados que no hayan sido verificados por el autor.

Todo con ayuda de una base de memoria que el autor fue construyendo poco a poco y unos documentos 
de notas que el autor fue tomando durante el desarrollo del proyecto con todo detalle y del análisis directo del código del proyecto.

\begin{itemize}
  \item El hallazgo más relevante de todos los detectados durante esta
        redacción es la distinción entre HOTA(0) y HOTA integrada
        : tanto las notas del
        autor como el propio script \texttt{summarize\_hyperparam.py}
        habían usado sistemáticamente la columna \texttt{HOTA(0)} de
        \texttt{TrackEval} bajo el nombre \enquote{HOTA}, lo que llevaba a una
        conclusión que se
        invierte con claridad al recalcular la comparativa con la métrica
        HOTA integrada estándar. Este hallazgo
        no estaba en las notas ni memoria originales: surgió de contrastar el código
        de evaluación con la documentación oficial de \texttt{TrackEval}
        al consultar con Claude la revisión de la parte de métricas en la documentación.
\end{itemize}

\section{Límites del uso de inteligencia artificial}
\label{sec:ia-limites}

Todo el contenido generado con Claude Code ha sido revisado y dirigido por
el autor en cada paso, no aceptado de forma automática:

\begin{itemize}
  \item Los cambios de alcance o de diseño con más de una alternativa
        razonable, por ejemplo, la decisión de generar metadatos de
        OC-SORT en vez de recortar el selector de algoritmo de la
        aplicación web, se plantearon
        explícitamente al autor como una decisión suya, no se resolvieron
        de forma autónoma por la herramienta.
  \item Todos los datos numéricos, tablas y resultados presentados en esta
        memoria proceden de ficheros realmente generados por el pipeline
        del proyecto (\texttt{evaluation\_train/}, \texttt{evaluation\_test/},
        \texttt{evaluation\_hyperparam/}, \texttt{metadata\_output/}), no de
        estimaciones ni de contenido inventado por la IA.
  \item Todo el código generado se ha ejecutado y su salida se ha
        inspeccionado antes de incorporarla a la memoria o de darla por
        válida (sección~\ref{sec:ia-desarrollo}); ningún resultado
        cuantitativo de este documento se ha aceptado sin volver a
        comprobar que procede realmente de la ejecución del sistema
        descrito.
\end{itemize}

En conjunto, el criterio seguido ha sido que Claude Code actúe como
herramienta de apoyo en la redacción y de desarrollo bajo dirección explícita del
autor en algunas decisiones, y que ninguna afirmación de esta
memoria se dé por buena sin haberla contrastado contra una
fuente verificable del proyecto.
